\chapter{2008年四川卷（理）}

\section{单选题}

\begin{problem}
设全集 \(U=\{1,2,3,4,5\}\)，\(A=\{1,2,3\}\)，\(B=\{2,3,4\}\)，则 \(\complement_U(A\cap B)=\)（\quad）

\choices
    {\(\{2,3\}\)}
    {\(\{1,4,5\}\)}
    {\(\{4,5\}\)}
    {\(\{1,5\}\)}
\end{problem}

\begin{answer}
B
\end{answer}

\begin{solution}
先求交集：\(A\cap B=\{2,3\}\). 因此 \(\complement_U(A\cap B)=U\mathbin{\backslash}\{2,3\}=\{1,4,5\}\)，故选 B.
\end{solution}

\begin{problem}
复数 \(2\i(1+\i)^2=\)（\quad）

\choices
    {\(-4\)}
    {\(4\)}
    {\(-4\i\)}
    {\(4\i\)}
\end{problem}

\begin{answer}
A
\end{answer}

\begin{solution}
由 \((1+\i)^2=1+2\i+\i^2=2\i\)，得 \(2\i(1+\i)^2=2\i\cdot2\i=4\i^2=-4\)，故选 A.
\end{solution}

\begin{problem}
\((\tan x+\cot x)\cos^2x=\)（\quad）

\choices
    {\(\tan x\)}
    {\(\sin x\)}
    {\(\cos x\)}
    {\(\cot x\)}
\end{problem}

\begin{answer}
D
\end{answer}

\begin{solution}
原式有意义时 \(\sin x\ne0\) 且 \(\cos x\ne0\). 于是 \((\tan x+\cot x)\cos^2x=\left(\dfrac{\sin x}{\cos x}+\dfrac{\cos x}{\sin x}\right)\cos^2x=\dfrac{\sin^2x+\cos^2x}{\sin x\cos x}\cos^2x=\dfrac{\cos x}{\sin x}=\cot x\)，故选 D.
\end{solution}

\begin{problem}
直线 \(y=3x\) 绕原点逆时针旋转 \(90^{\circ}\)，再向右平移 \(1\) 个单位，所得到的直线为（\quad）

\choices
    {\(y=-\dfrac13x+\dfrac13\)}
    {\(y=-\dfrac13x+1\)}
    {\(y=3x-3\)}
    {\(y=\dfrac13x+1\)}
\end{problem}

\begin{answer}
A
\end{answer}

\begin{solution}
直线 \(y=3x\) 绕原点逆时针旋转 \(90^{\circ}\) 后，与原直线垂直且过原点，故方程为 \(y=-\dfrac13x\). 向右平移 \(1\) 个单位后，把新横坐标 \(x\) 代回平移前的方程，得 \(y=-\dfrac13(x-1)=-\dfrac13x+\dfrac13\)，故选 A.
\end{solution}

\begin{problem}
若 \(0\leqslant\alpha\leqslant2\pi\)，\(\sin\alpha>\sqrt3\cos\alpha\)，则 \(\alpha\) 的取值范围是（\quad）

\choices
    {\(\left(\dfrac\pi3,\dfrac\pi2\right)\)}
    {\(\left(\dfrac\pi3,\pi\right)\)}
    {\(\left(\dfrac\pi3,\dfrac{4\pi}3\right)\)}
    {\(\left(\dfrac\pi3,\dfrac{3\pi}2\right)\)}
\end{problem}

\begin{answer}
C
\end{answer}

\begin{solution}
由辅助角公式，\(\sin\alpha-\sqrt3\cos\alpha=2\sin\left(\alpha-\dfrac\pi3\right)\)，所以题设不等式等价于 \(\sin\left(\alpha-\dfrac\pi3\right)>0\). 当 \(0\leqslant\alpha\leqslant2\pi\) 时，\(-\dfrac\pi3\leqslant\alpha-\dfrac\pi3\leqslant\dfrac{5\pi}3\)，正弦为正的部分为 \(0<\alpha-\dfrac\pi3<\pi\)，即 \(\dfrac\pi3<\alpha<\dfrac{4\pi}3\)，故选 C.
\end{solution}

\begin{problem}
从甲、乙等 \(10\) 个同学中挑选 \(4\) 名参加某项公益活动，要求甲、乙中至少有 \(1\) 人参加，则不同的挑选方法共有（\quad）

\choices
    {\(70\) 种}
    {\(112\) 种}
    {\(140\) 种}
    {\(168\) 种}
\end{problem}

\begin{answer}
C
\end{answer}

\begin{solution}
从 \(10\) 个同学中任取 \(4\) 人共有 \(\mathrm C_{10}^{4}\) 种. 若甲、乙均未参加，则只能从其余 \(8\) 人中取 \(4\) 人，共有 \(\mathrm C_{8}^{4}\) 种. 因此符合条件的方法数为 \(\mathrm C_{10}^{4}-\mathrm C_{8}^{4}=210-70=140\)，故选 C.
\end{solution}

\begin{problem}
已知等比数列 \(\{a_n\}\) 中 \(a_2=1\)，则其前 \(3\) 项的和 \(S_3\) 的取值范围是（\quad）

\choices
    {\((-\infty,-1]\)}
    {\((-\infty,0)\cup(1,+\infty)\)}
    {\([3,+\infty)\)}
    {\((-\infty,-1]\cup[3,+\infty)\)}
\end{problem}

\begin{answer}
D
\end{answer}

\begin{solution}
设等比数列的公比为 \(q\). 由 \(a_2=1\) 得 \(a_1=\dfrac1q\)，且 \(q\ne0\)，所以 \(S_3=\dfrac1q+1+q=1+q+\dfrac1q\). 当 \(q>0\) 时，\(q+\dfrac1q\geqslant2\)；当 \(q<0\) 时，\(q+\dfrac1q\leqslant-2\). 因此 \(S_3\in(-\infty,-1]\cup[3,+\infty)\)，故选 D.
\end{solution}

\begin{problem}
设 \(M\)，\(N\) 是球心 \(O\) 的半径 \(OP\) 上的两点，且 \(NP=MN=OM\)，分别过 \(N\)，\(M\)，\(O\) 作垂直于 \(OP\) 的平面，截球面得三个圆，则这三个圆的面积之比为（\quad）

\choices
    {\(3:5:6\)}
    {\(3:6:8\)}
    {\(5:7:9\)}
    {\(5:8:9\)}
\end{problem}

\begin{answer}
D
\end{answer}

\begin{solution}
设球半径 \(OP=R\)，取 \(OP\) 为数轴正向，设 \(OM=x>0\)，\(ON=t\). 由 \(NP=OM\) 和 \(MN=OM\)，得 \(R-t=x\)，\(\lvert t-x\rvert=x\)，所以 \(t=0\) 或 \(t=2x\). 若 \(t=0\)，则 \(N=O\)，且 \(R=x\)，从而 \(M=P\)，过 \(N\) 与过 \(O\) 的截面平面重合，不符合题设得到三个截面圆的情形，故 \(t=2x\). 因而 \(ON=2x\)，\(OP=3x\)，即 \(O\)，\(M\)，\(N\)，\(P\) 依次位于半径 \(OP\) 上.

球被过 \(N\)、\(M\)、\(O\) 且垂直于 \(OP\) 的平面截得的三个圆，其圆心到球心的距离分别为 \(2x\)、\(x\)、\(0\). 所以三个截面圆半径的平方之比为 \((3x)^2-(2x)^2:(3x)^2-x^2:(3x)^2=5:8:9\). 圆面积与半径平方成正比，故选 D.
\end{solution}

\begin{problem}
设直线 \(l\subset\) 平面 \(\alpha\)，过平面 \(\alpha\) 外一点 \(A\) 与 \(l\)，\(\alpha\) 都成 \(30^{\circ}\) 角的直线有且只有（\quad）

\choices
    {\(1\) 条}
    {\(2\) 条}
    {\(3\) 条}
    {\(4\) 条}
\end{problem}

\begin{answer}
B
\end{answer}

\begin{solution}
设所求直线为 \(m\)，其在平面 \(\alpha\) 上的正投影为 \(m'\)，设 \(m'\) 与 \(l\) 的锐角为 \(\theta\). 因为 \(m\) 与平面 \(\alpha\) 的角为 \(30^{\circ}\)，而 \(l\subset\alpha\)，由方向向量的投影关系有 \(\cos30^{\circ}=\cos30^{\circ}\cos\theta\)，从而 \(\cos\theta=1\)，即 \(m'\parallel l\).

设 \(\bs u\) 是 \(l\) 的单位方向向量，设 \(\bs n\) 是由平面 \(\alpha\) 指向 \(A\) 的单位法向量. 满足条件的直线方向可取为 \(\cos30^{\circ}\bs u-\sin30^{\circ}\bs n\) 或 \(-\cos30^{\circ}\bs u-\sin30^{\circ}\bs n\). 这两个方向不共线，且除这两种相反的投影方向外没有其他可能，所以共有 \(2\) 条，故选 B.
\end{solution}

\begin{problem}
设 \(f(x)=\sin(\omega x+\varphi)\)，其中 \(\omega>0\)，则 \(f(x)\) 是偶函数的充要条件是（\quad）

\choices
    {\(f(0)=1\)}
    {\(f(0)=0\)}
    {\(f'(0)=1\)}
    {\(f'(0)=0\)}
\end{problem}

\begin{answer}
D
\end{answer}

\begin{solution}
利用和角公式，\(f(x)=\sin\varphi\cos\omega x+\cos\varphi\sin\omega x\). 其中第一项是偶函数，第二项是奇函数，因此 \(f(x)\) 为偶函数当且仅当 \(\cos\varphi=0\). 又 \(f'(0)=\omega\cos\varphi\)，且 \(\omega>0\)，故 \(f(x)\) 为偶函数当且仅当 \(f'(0)=0\)，选 D.
\end{solution}

\begin{problem}
设定义在 \(\R\) 上的函数 \(f(x)\) 满足 \(f(x)\cdot f(x+2)=13\)，若 \(f(1)=2\)，则 \(f(99)=\)（\quad）

\choices
    {\(13\)}
    {\(2\)}
    {\(\dfrac{13}{2}\)}
    {\(\dfrac{2}{13}\)}
\end{problem}

\begin{answer}
C
\end{answer}

\begin{solution}
由 \(f(1)=2\) 及 \(f(x)f(x+2)=13\)，得 \(f(3)=\dfrac{13}{f(1)}=\dfrac{13}{2}\)，继续得到 \(f(5)=\dfrac{13}{f(3)}=2\). 因此在奇数点上，函数值按 \(2\)，\(\dfrac{13}{2}\) 交替出现. 写成 \(f(1+2k)\) 时，\(k\) 为偶数取 \(2\)，\(k\) 为奇数取 \(\dfrac{13}{2}\). 由于 \(99=1+2\times49\)，\(49\) 为奇数，所以 \(f(99)=\dfrac{13}{2}\)，故选 C.
\end{solution}

\begin{problem}
已知抛物线 \(C:y^2=8x\) 的焦点为 \(F\)，准线与 \(x\) 轴的交点为 \(K\)，点 \(A\) 在 \(C\) 上，且 \(\lvert AK\rvert=\sqrt2\lvert AF\rvert\)，则 \(\triangle AFK\) 的面积为（\quad）

\choices
    {\(4\)}
    {\(8\)}
    {\(16\)}
    {\(32\)}
\end{problem}

\begin{answer}
B
\end{answer}

\begin{solution}
将 \(y^2=8x\) 与标准式 \(y^2=4px\) 对比，得 \(p=2\)，所以焦点 \(F=(2,0)\)，准线为 \(x=-2\)，其与 \(x\) 轴交点为 \(K=(-2,0)\). 设 \(A=(x_0,y_0)\)，则由抛物线的焦点与准线定义，\(AF=x_0+2\)，而 \(AK^2=(x_0+2)^2+y_0^2\). 条件 \(AK=\sqrt2AF\) 给出 \((x_0+2)^2+y_0^2=2(x_0+2)^2\)，即 \(y_0^2=(x_0+2)^2\). 结合 \(y_0^2=8x_0\)，得 \(8x_0=(x_0+2)^2\)，从而 \((x_0-2)^2=0\)，故 \(x_0=2\)，\(y_0=\pm4\).

三角形 \(AFK\) 的底 \(FK=4\)，高为 \(\lvert y_0\rvert=4\)，所以面积为 \(\dfrac12\times4\times4=8\)，故选 B.
\end{solution}

\section{填空题}

\begin{problem}
\((1+2x)^3(1-x)^4\) 展开式中 \(x^2\) 的系数为\fillinblank{}.
\end{problem}

\begin{answer}
\(-6\)
\end{answer}

\begin{solution}
展开前两项，\((1+2x)^3=1+6x+12x^2+8x^3\)，\((1-x)^4=1-4x+6x^2-4x^3+x^4\). 要得到 \(x^2\) 项，只有三种相乘方式，因此其系数为 \(1\cdot6+6\cdot(-4)+12\cdot1=6-24+12=-6\).
\end{solution}

\begin{problem}
已知直线 \(l:x-y+4=0\) 与圆 \(C:(x-1)^2+(y-1)^2=2\)，则 \(C\) 上各点到 \(l\) 的距离的最小值为\fillinblank{}.
\end{problem}

\begin{answer}
\(\sqrt{2}\)
\end{answer}

\begin{solution}
圆心为 \(O=(1,1)\)，半径为 \(r=\sqrt2\). 圆心到直线 \(l\) 的距离为 \(\dfrac{\lvert1-1+4\rvert}{\sqrt{1^2+(-1)^2}}=2\sqrt2\). 由于 \(2\sqrt2>\sqrt2\)，直线与圆不相交，圆上各点到直线的最小距离等于圆心距减去半径，即 \(2\sqrt2-\sqrt2=\sqrt2\).
\end{solution}

\begin{problem}
已知正四棱柱的对角线的长为 \(\sqrt6\)，且对角线与底面所成角的余弦值为 \(\dfrac{\sqrt3}{3}\)，则该正四棱柱的体积等于\fillinblank{}.
\end{problem}

\begin{answer}
2
\end{answer}

\begin{solution}
设正四棱柱底面边长为 \(a\)，高为 \(h\)，对角线长为 \(d=\sqrt6\). 对角线在底面上的投影是正方形的对角线，长度为 \(a\sqrt2\). 因此 \(\cos\theta=\dfrac{a\sqrt2}{\sqrt6}=\dfrac{\sqrt3}{3}\)，解得 \(a=1\). 又 \(d^2=2a^2+h^2\)，所以 \(6=2+h^2\)，得 \(h=2\). 故体积为 \(a^2h=1^2\cdot2=2\).
\end{solution}

\begin{problem}
设等差数列 \(\{a_n\}\) 的前 \(n\) 项和为 \(S_n\)，若 \(S_4\geqslant10\)，\(S_5\leqslant15\)，则 \(a_4\) 的最大值为\fillinblank{}.
\end{problem}

\begin{answer}
4
\end{answer}

\begin{solution}
设等差数列的首项为 \(a_1\)，公差为 \(d\). 由前 \(n\) 项和公式，\(S_4=4a_1+6d\)，\(S_5=5a_1+10d\)，故
\(2a_1+3d\geqslant5\)，\(a_1+2d\leqslant3\). 前一个不等式给出 \(a_1\geqslant\dfrac{5-3d}{2}\)，与后一个不等式结合可得 \(\dfrac{5-3d}{2}\leqslant3-2d\)，从而 \(d\leqslant1\). 于是 \(a_4=a_1+3d\leqslant3-2d+3d=3+d\leqslant4\). 取 \(d=1\)，\(a_1=1\) 时，\(S_4=10\)，\(S_5=15\)，且 \(a_4=4\)，所以最大值为 \(4\).
\end{solution}

\section{解答题}

\begin{problem}
求函数 \(y=7-4\sin x\cos x+4\cos^2x-4\cos^4x\) 的最大值与最小值.
\end{problem}

\begin{solution}
利用 \(2\sin x\cos x=\sin2x\) 以及 \(\sin^2x\cos^2x=\dfrac14\sin^22x\)，有 \(4\cos^2x-4\cos^4x=4\sin^2x\cos^2x=\sin^22x\). 因此
\(y=7-2\sin2x+\sin^22x=(\sin2x-1)^2+6\). 令 \(u=\sin2x\)，则 \(-1\leqslant u\leqslant1\)，从而 \(0\leqslant(u-1)^2\leqslant4\). 所以 \(6\leqslant y\leqslant10\)，最大值为 \(10\)，最小值为 \(6\).
\end{solution}

\begin{problem}
设进入某商场的每一位顾客购买甲种商品的概率为 \(0.5\)，购买乙种商品的概率为 \(0.6\)，且购买甲种商品与购买乙种商品相互独立，各顾客之间购买商品也是相互独立的.

\begin{enumerate}
\item 求进入商场的 \(1\) 位顾客购买甲、乙两种商品中的一种的概率；
\item 求进入商场的 \(1\) 位顾客至少购买甲、乙两种商品中的一种的概率；
\item 记 \(\xi\) 表示进入商场的 \(3\) 位顾客中至少购买甲、乙两种商品中的一种的人数，求 \(\xi\) 的分布列及期望.
\end{enumerate}
\end{problem}

\begin{solution}
\begin{enumerate}
\item 设顾客购买甲、乙商品分别为事件 \(A\)、\(B\)，则 \(P(A)=0.5\)，\(P(B)=0.6\). 购买其中一种表示恰购买一种，因此概率为 \(P(A\overline B)+P(\overline A B)=0.5\times0.4+0.5\times0.6=0.5\).
\item 至少购买一种商品的对立事件是两种商品都不购买，故概率为 \(1-P(\overline A\overline B)=1-0.5\times0.4=0.8\).
\item 对每位顾客，购买至少一种商品的概率为 \(p=0.8\)，且顾客之间相互独立，所以 \(\xi\sim B(3,0.8)\). 对 \(k=0\)，\(1\)，\(2\)，\(3\)，\(P(\xi=k)=\mathrm C_3^k(0.8)^k(0.2)^{3-k}\)，因而
\begin{center}
\begin{tabular}{c|cccc}
\(k\) & \(0\) & \(1\) & \(2\) & \(3\) \\
\hline
\(P(\xi=k)\) & \(0.008\) & \(0.096\) & \(0.384\) & \(0.512\)
\end{tabular}
\end{center}
最后 \(E(\xi)=np=3\times0.8=2.4\).
\end{enumerate}
\end{solution}

\begin{problem}
如图，平面 \(ABEF\perp\) 平面 \(ABCD\)，四边形 \(ABEF\) 与 \(ABCD\) 都是直角梯形，\(\angle BAD=\angle FAB=90^{\circ}\)，\(BC\parallel\dfrac12AD\)，\(BE\parallel\dfrac12AF\).

\begin{enumerate}
\item 证明：\(C\)，\(D\)，\(F\)，\(E\) 四点共面；
\item 设 \(AB=BC=BE\)，求二面角 \(A-ED-B\) 的大小.
\end{enumerate}

\begin{center}
\bitmapfigure[width=0.60\linewidth]{img/2008/sichuan_ordinary_science/q19_fig1.png}
\end{center}
\end{problem}

\begin{solution}
\begin{enumerate}
\item 由题设的比例关系可写成 \(\overrightarrow{BC}=\dfrac12\overrightarrow{AD}\)，\(\overrightarrow{BE}=\dfrac12\overrightarrow{AF}\). 于是 \(\overrightarrow{CE}=\overrightarrow{BE}-\overrightarrow{BC}=\dfrac12(\overrightarrow{AF}-\overrightarrow{AD})=\dfrac12\overrightarrow{DF}\)，所以 \(CE\parallel DF\)，四点 \(C\)，\(D\)，\(F\)，\(E\) 共面.
\item 设 \(AB=BC=BE=1\)，则 \(AD=AF=2\). 由两平面垂直且交线为 \(AB\)，并结合 \(AD\perp AB\)、\(AF\perp AB\)，可建立以 \(A\) 为原点、\(AB\)、\(AD\)、\(AF\) 为三个互相垂直坐标轴的直角坐标系. 于是 \(A=(0,0,0)\)，\(B=(1,0,0)\)，\(D=(0,2,0)\)，\(F=(0,0,2)\)，\(C=(1,1,0)\)，\(E=(1,0,1)\).

平面 \(AED\) 的一个法向量为 \(\bs n_1=(1,0,-1)\)，平面 \(BED\) 的一个法向量为 \(\bs n_2=(2,1,0)\). 设二面角 \(A-ED-B\) 的平面角为 \(\theta\)，则 \(\cos\theta=\dfrac{\bs n_1\cdot\bs n_2}{\lvert\bs n_1\rvert\lvert\bs n_2\rvert}=\dfrac{2}{\sqrt{10}}\). 因此 \(\tan\theta=\dfrac{\sqrt6}{2}\)，故二面角的大小为 \(\arctan\dfrac{\sqrt6}{2}\).
\end{enumerate}
\end{solution}

\begin{problem}
设数列 \(\{a_n\}\) 的前 \(n\) 项和为 \(S_n\)，已知 \(ba_n-2^n=(b-1)S_n\).

\begin{enumerate}
\item 证明：当 \(b=2\) 时，\(\{a_n-n\cdot2^{n-1}\}\) 是等比数列；
\item 求 \(\{a_n\}\) 的通项公式.
\end{enumerate}
\end{problem}

\begin{solution}
由 \(n=1\) 时的条件得 \(ba_1-2=(b-1)a_1\)，所以 \(a_1=2\). 对 \(n\geqslant2\)，将第 \(n\) 项和第 \(n-1\) 项的条件相减，得 \(a_n=ba_{n-1}+2^{n-1}\).

\begin{enumerate}
\item 当 \(b=2\) 时，令 \(u_n=a_n-n\cdot2^{n-1}\). 由递推式，
\(u_n=a_n-n\cdot2^{n-1}=2a_{n-1}+2^{n-1}-n\cdot2^{n-1}=2\left(a_{n-1}-(n-1)\cdot2^{n-2}\right)=2u_{n-1}\). 又 \(u_1=a_1-1=1\)，所以 \(\{u_n\}\) 是首项为 \(1\)、公比为 \(2\) 的等比数列，即 \(\{a_n-n\cdot2^{n-1}\}\) 是等比数列.
\item 当 \(b=2\) 时，由上一问 \(u_n=2^{n-1}\)，所以 \(a_n=(n+1)2^{n-1}\). 当 \(b\ne2\) 时，递推式迭代得
\(a_n=2b^{n-1}+\sum_{k=2}^{n}b^{n-k}2^{k-1}=2b^{n-1}+\dfrac{2^n-2b^{n-1}}{2-b}=\dfrac{2^n+2(1-b)b^{n-1}}{2-b}\). 因此
\[
a_n=
\begin{cases}
(n+1)2^{n-1}, & b=2,\\
\dfrac{2^n}{2-b}+\dfrac{2(1-b)}{2-b}b^{n-1}, & b\ne2.
\end{cases}
\]
\end{enumerate}
\end{solution}

\begin{problem}
设椭圆 \(\dfrac{x^2}{a^2}+\dfrac{y^2}{b^2}=1\)（\(a>b>0\)）的左、右焦点分别为 \(F_1\)，\(F_2\)，离心率 \(e=\dfrac{\sqrt2}{2}\)，右准线为 \(l\)，\(M\)，\(N\) 是 \(l\) 上的两个动点，\(\overrightarrow{F_1M}\cdot\overrightarrow{F_2N}=0\).

\begin{enumerate}
\item 若 \(\lvert\overrightarrow{F_1M}\rvert=\lvert\overrightarrow{F_2N}\rvert=2\sqrt5\)，求 \(a\)，\(b\) 的值；
\item 证明：当 \(\lvert MN\rvert\) 取最小值时，\(\overrightarrow{F_1M}+\overrightarrow{F_2N}\) 与 \(\overrightarrow{F_1F_2}\) 共线.
\end{enumerate}

\FigureLayoutDeclare{img/2008/sichuan_ordinary_science/q21_fig1.png}{11.0cm}{30bp 36bp 17bp 40bp}
\begin{center}
\bitmapfigure[width=0.65\linewidth]{img/2008/sichuan_ordinary_science/q21_fig1.png}
\end{center}
\end{problem}

\begin{solution}
设椭圆焦距的一半为 \(c\)，右准线横坐标为 \(d\)，则 \(c=ae=\dfrac{a\sqrt2}{2}\)，\(d=\dfrac ae=a\sqrt2\). 取 \(M=(d,m)\)，\(N=(d,n)\)，于是 \(F_1=(-c,0)\)，\(F_2=(c,0)\)，从而 \(d+c=\dfrac{3a\sqrt2}{2}\)，\(d-c=\dfrac{a\sqrt2}{2}\).

\begin{enumerate}
\item 由两条线段长度条件，得 \(m^2=20-\dfrac{9a^2}{2}\)，\(n^2=20-\dfrac{a^2}{2}\). 又由垂直条件，
\(0=\overrightarrow{F_1M}\cdot\overrightarrow{F_2N}=(d+c)(d-c)+mn=\dfrac{3a^2}{2}+mn\)，所以 \(mn=-\dfrac{3a^2}{2}\). 两边平方并代入 \(m^2\)，\(n^2\)，得
\(\left(20-\dfrac{9a^2}{2}\right)\left(20-\dfrac{a^2}{2}\right)=\dfrac{9a^4}{4}\)，化简为 \(400-100a^2=0\)，故 \(a=2\). 再由 \(b^2=a^2(1-e^2)\)，得 \(b^2=4\left(1-\dfrac12\right)=2\)，所以 \(b=\sqrt2\).
\item 由垂直条件的一般形式，\(mn=-(d^2-c^2)\). 令 \(K=d^2-c^2>0\)，则
\(\lvert MN\rvert^2=(m-n)^2=m^2+n^2-2mn=m^2+n^2+2K\geqslant4K\)，其中等号成立当且仅当 \(m^2=n^2\). 结合 \(mn=-K<0\)，最小值取得时 \(n=-m\).

此时 \(\overrightarrow{F_1M}+\overrightarrow{F_2N}=(d+c,m)+(d-c,n)=(2d,m+n)=(2d,0)\)，而 \(\overrightarrow{F_1F_2}=(2c,0)\)，所以两向量共线.
\end{enumerate}
\end{solution}

\begin{problem}
已知 \(x=3\) 是函数 \(f(x)=a\ln(1+x)+x^2-10x\) 的一个极值点.

\begin{enumerate}
\item 求 \(a\)；
\item 求函数 \(f(x)\) 的单调区间；
\item 若直线 \(y=b\) 与函数 \(y=f(x)\) 的图象有 \(3\) 个交点，求 \(b\) 的取值范围.
\end{enumerate}
\end{problem}

\begin{solution}
函数定义域为 \(x>-1\)，且 \(f'(x)=\dfrac{a}{1+x}+2x-10\).

\begin{enumerate}
\item \(x=3\) 是极值点，故 \(f'(3)=0\)，即 \(\dfrac a4+6-10=0\)，解得 \(a=16\).
\item 此时 \(f'(x)=\dfrac{16}{1+x}+2x-10=\dfrac{2(x-1)(x-3)}{x+1}\). 在定义域内 \(x+1>0\)，所以当 \(-1<x<1\) 或 \(x>3\) 时 \(f'(x)>0\)，当 \(1<x<3\) 时 \(f'(x)<0\). 因此增区间为 \((-1,1)\) 和 \((3,+\infty)\)，减区间为 \((1,3)\).
\item 由上面的单调性，\(f(x)\) 在 \((-1,1)\) 上从 \(-\infty\) 增至 \(f(1)=16\ln2-9\)，在 \((1,3)\) 上减至 \(f(3)=32\ln2-21\)，在 \((3,+\infty)\) 上再增至 \(+\infty\). 水平直线与三段图象各有一个交点，当且仅当 \(f(3)<b<f(1)\). 在端点处会与极值点重合，只有两个不同交点，所以 \(b\in(32\ln2-21,16\ln2-9)\).
\end{enumerate}
\end{solution}
